\subsection{EXAMPLE: POLYNOMIALS WITH INTEGER VALUES}

Let V be a finite dimensional \(\mathbf{Q}\)-vector space, \(V^{*}\) its dual, \(\mathcal{V}\) a lattice in V, \(\mathcal{V}^{*}\) the dual \(\mathbf{Z}\)-module of \(\mathcal{V}\), which can be identified canonically with a lattice in \(V^{*}\) , \(\mathbf{S}(V)\) the symmetric algebra of V, and

\[
\lambda : \mathbf {S} (\mathrm{V}) \to \mathrm{A} (\mathrm{V} ^ {*})
\]

the canonical bijection from \(\mathbf{S}(\mathrm{V})\) to the algebra of polynomial functions on \(\mathrm{V}^*\) (Algebra, Chap. IV, §5, no. 11, Remark 1). If we identify \(\mathrm{A}(\mathrm{V}^* \times \mathrm{V}^*)\) with \(\mathrm{A}(\mathrm{V}^*) \otimes_{\mathbf{Q}} \mathrm{A}(\mathrm{V}^*)\) , then \(\lambda\) transforms the coproduct of \(\mathbf{S}(\mathrm{V})\) into the map \(\mathrm{A}(\mathrm{V}^*) \to \mathrm{A}(\mathrm{V}^* \times \mathrm{V}^*)\) which associates to the polynomial function \(\varphi\) on \(\mathrm{V}^*\) the polynomial function

\[
(x, y) \mapsto \varphi (x + y)
\]

on \(\mathrm{V}^* \times \mathrm{V}^*\) (Algebra, Chap. IV, §5, no. 11, Remark 2).

Denote by \(\binom{\mathcal{V}}{\mathbf{Z}}\) the subset of \(\mathbf{S}(\mathrm{V})\) consisting of the elements which correspond to polynomial maps from \(\mathrm{V}^*\) to \(\mathbf{Q}\) that take integer values on \(\mathcal{V}^*\) .

PROPOSITION 2. (i) \(\begin{pmatrix} \mathcal{V} \\ \mathbf{Z} \end{pmatrix}\) is a biorder in the bigebra \(\mathbf{S}(\mathrm{V})\) , and \(\begin{pmatrix} \mathcal{V} \\ \mathbf{Z} \end{pmatrix} \cap \mathrm{V} = \mathcal{V}\) .

\begin{enumerate}
\def\labelenumi{(\roman{enumi})}
\setcounter{enumi}{1}
\item
  The \(\mathbf{Z}\) -algebra \(\left( \begin{array}{c}\mathcal{V}\\ \mathbf{Z} \end{array} \right)\) is generated by the \(\left( \begin{array}{c}h\\ n \end{array} \right)\) for \(h\in \mathcal{V},n\in \mathbf{N}\) .
\item
  If \((h_1, \ldots, h_r)\) is a basis of \(\mathcal{V}\) , the elements
\end{enumerate}

\[
\binom{h}{\mathbf {n}} = \binom{h _ {1}}{n _ {1}} \dots \binom{h _ {r}}{n _ {r}},
\]

where \(\mathbf{n} = (n_1, \ldots, n_r)\) belongs to \(\mathbf{N}^r\) , form a basis of the \(\mathbf{Z}\) -module \(\binom{\mathcal{V}}{\mathbf{Z}}\) .

For \(m \in \mathbf{N}\) , put \(\mathbf{S}_{m}(\mathrm{V}) = \sum_{i \leq m} \mathbf{S}^{i}(\mathrm{V}), \mathbf{S}_{m}(\mathcal{V}) = \sum_{i \leq m} \mathbf{S}^{i}(\mathcal{V})\) . By Algebra, Chap. IV, §5, no. 9, Prop. 15 and Remark,

\[
\mathbf {S} _ {m} (\mathcal {V}) \subset \mathbf {S} _ {m} (\mathrm{V}) \cap \binom{\mathcal {V}}{\mathbf {Z}} \subset \frac {1}{m !} \mathbf {S} _ {m} (\mathcal {V})
\]

so \(\binom{\mathcal{V}}{\mathbf{Z}}\cap\mathrm{V}=\mathcal{V}\) . Since \(\mathbf{S}_{m}(\mathcal{V})\) is a lattice in \(\mathbf{S}_{m}(\mathrm{V})\) , \(\mathbf{S}_{m}(\mathrm{V})\cap\binom{\mathcal{V}}{\mathbf{Z}}\) is also a lattice in \(\mathbf{S}_{m}(\mathrm{V})\) . On the other hand, \(\mathbf{S}_{m}(\mathrm{V})\cap\binom{\mathcal{V}}{\mathbf{Z}}\) is a direct factor of \(\mathbf{S}_{m+1}(\mathrm{V})\cap\binom{\mathcal{V}}{\mathbf{Z}}\) (since the quotient is torsion-free), hence it admits a complement which is a free \(\mathbf{Z}\)-module. It follows that \(\binom{\mathcal{V}}{\mathbf{Z}}\) is a free \(\mathbf{Z}\)-module. It is clear that this is a unital order in the algebra \(\mathbf{S}(\mathrm{V})\) . Let \((u_{n})_{n\in\mathbf{N}}\) be a basis of the \(\mathbf{Z}\)-module \(\binom{\mathcal{V}}{\mathbf{Z}}\) . This is also a basis of the \(\mathbf{Q}\)-module \(\mathbf{S}(\mathrm{V})\) and, for all

\[
\varphi \in \mathbf {S} (\mathrm{V} \times \mathrm{V}) = \mathbf {S} (\mathrm{V}) \otimes_ {\mathbf {Q}} \mathbf {S} (\mathrm{V}),
\]

there exists a unique sequence \((v_{n})\) of elements of S(V) such that \(\varphi=\sum u_{n}\otimes v_{n}\) . As above, identify S(V) with A(V*) and S(V) \(\otimes\) S(V) with A(V* × V*). If \(\varphi\in\binom{\mathcal{V}\times\mathcal{V}}{\mathbf{Z}}\) , the polynomial function \(x\mapsto\varphi(x,y)\) belongs to \(\binom{\mathcal{V}}{\mathbf{Z}}\) for all \(y\in\mathcal{V}^{*}\) . It follows that \(v_{n}(y)\in\mathbf{Z}\) for all n and all \(y\in\mathcal{V}^{*}\) , in other words that \(v_{n}\in\binom{\mathcal{V}}{\mathbf{Z}}\) . This proves that the coproduct maps \(\binom{\mathcal{V}}{\mathbf{Z}}\) to \(\binom{\mathcal{V}}{\mathbf{Z}}\otimes_{\mathbf{Z}}\binom{\mathcal{V}}{\mathbf{Z}}\) . If \(h\in\mathcal{V}\) and \(n\in\mathbf{N}\) , then \(\binom{h}{n}\) maps \(u\in\mathcal{V}^{*}\) to the integer \(\binom{u(h)}{n}\) , so \(\binom{h}{n}\in\binom{\mathcal{V}}{\mathbf{Z}}\) . Assertion (iii) is now obtained by applying Th. 1 to the commutative Lie algebra V, and (ii) follows.

COROLLARY. Let \(\mathrm{X}\) be an indeterminate. The polynomials \(\binom{\mathrm{X}}{n}\) , where \(n \in \mathbf{N}\) , form a basis of the \(\mathbf{Z}\) -module consisting of the polynomials \(\mathrm{P} \in k[\mathrm{X}]\) such that \(\mathrm{P}(\mathbf{Z}) \subset \mathbf{Z}\) .

If \(\mathrm{P}(\mathbf{Z}) \subset \mathbf{Z}\) , the Lagrange interpolation formula (Algebra, Chap. IV, §2, no. 1, Prop. 6) shows that the coefficients of P belong to \(\mathbf{Q}\) ; thus, we can assume that \(k = \mathbf{Q}\) and apply Prop. 2 with \(V = \mathbf{Q}\) , \(\mathcal{V} = \mathbf{Z}\) .

\subsection{SOME FORMULAS}

In this number, A denotes a unital associative algebra. If \(x \in A\) , we write ad x instead of \(ad_{A}x\) .

Lemma 2. If \(x, y \in \mathbf{A}\) and \(n \in \mathbf{N}\) ,

\[
\frac {(\operatorname{ad} x) ^ {n}}{n !} y = \sum_ {p + q = n} (- 1) ^ {q} \frac {x ^ {p}}{p !} y \frac {x ^ {q}}{q !} = \sum_ {p + q = n} (- 1) ^ {q} x ^ {(p)} y x ^ {(q)}.\tag{11}
\]

Indeed, if we denote by \(L_{x}\) and \(R_{x}\) the maps \(z \mapsto xz\) and \(z \mapsto zx\) from A to A, we have, since \(L_{x}\) and \(R_{x}\) commute,

\[
\frac {1}{n !} (\mathrm{ad} x) ^ {n} = \frac {1}{n !} (\mathrm{L} _ {x} - \mathrm{R} _ {x}) ^ {n} = \sum_ {p + q = n} (- 1) ^ {q} \frac {1}{p !} \mathrm{L} _ {x} ^ {p} \frac {1}{q !} \mathrm{R} _ {x} ^ {q}.
\]

Lemma 3. Let \(x, h \in \mathbf{A}\) and \(\lambda \in k\) be such that (ad \(h)x = \lambda x\) . For all \(n \in \mathbf{N}\) , and all \(\mathrm{P} \in k[\mathrm{X}]\) , we have

\[
\mathrm{P} (h) x ^ {(n)} = x ^ {(n)} \mathrm{P} (h + n \lambda).\tag{12}
\]

Since \(\operatorname{ad} h\) is a derivation of \(A\) and since \((\operatorname{ad} h)x\) commutes with \(x\) , we have

\[
(\operatorname{ad} h) x ^ {n} = n x ^ {n - 1} ((\operatorname{ad} h) x) = n \lambda x ^ {n},\tag{13}
\]

SO

\[
(\operatorname{ad} h) x ^ {(n)} = n \lambda x ^ {(n)}.
\]

Thus, formula (12) follows from the special case

\[
\mathrm{P} (h) x = x \mathrm{P} (h + \lambda)\tag{14}
\]

by replacing \(x\) by \(x^{(n)}\) and \(\lambda\) by \(n\lambda\) . It suffices to prove (14) when \(\mathrm{P} = \mathrm{X}^m\) , by induction on \(m\) . It is clear when \(m = 0, 1\) . If (14) is true for \(\mathrm{P} = \mathrm{X}^m\) , then

\[
h ^ {m + 1} x = h. h ^ {m} x = h x (h + \lambda) ^ {m} = x (h + \lambda) ^ {m + 1}
\]

which proves (12).

Lemma 4. Let \(x, y, h \in \mathbf{A}\) be such that

\[
[ y, x ] = h, \quad [ h, x ] = 2 x, \quad [ h, y ] = - 2 y.\tag{15}
\]

\begin{enumerate}
\def\labelenumi{(\roman{enumi})}
\tightlist
\item
  For \(m, n \in \mathbf{N}\) , we have
\end{enumerate}

\[
x ^ {(n)} y ^ {(m)} = \sum_ {p \geq 0} y ^ {(m - p)} \binom{m + n - p - 1 - h}{p} x ^ {(n - p)}.\tag{16}
\]

\begin{enumerate}
\def\labelenumi{(\roman{enumi})}
\setcounter{enumi}{1}
\tightlist
\item
  Let \(\mathbf{A}'\) be the \(\mathbf{Z}\) -subalgebra of \(\mathbf{A}\) generated by the \(x^{(m)}\) and the \(y^{(m)}\) for \(m \in \mathbf{N}\) . Then \(\binom{h}{n} \in \mathbf{A}'\) for all \(n \in \mathbf{N}\) .
\end{enumerate}

Formula (16) can be written in the equivalent form

\[
(\operatorname{ad} x ^ {(n)}) y ^ {(m)} = \sum_ {p \geq 1} y ^ {(m - p)} \binom {m + n - p - 1 - h} {p} x ^ {(n - p)}.\tag{$17_{m}$}
\]

This is trivial for \(m = 0\) . We argue by induction on \(m\) . From \((17_{m})\) , we obtain

\[
\begin{array}{l} (m + 1) (\operatorname{ad} x ^ {(n)}) y ^ {(m + 1)} = (\operatorname{ad} x ^ {(n)}) y ^ {(m)}. y + y ^ {(m)}. (\operatorname{ad} x ^ {(n)}) y \\ = \sum_ {p \geq 1} y ^ {(m - p)} \binom {m + n - p - 1 - h} {p} x ^ {(n - p)} y + y ^ {(m)} (n - 1 - h) x ^ {(n - 1)} \end{array} \tag {18}
\]

(§1, no. 1, Lemma 1). Now, applying the same lemma, and then Lemma 3, we have

\[
\begin{array}{l} \binom {m + n - p - 1 - h} {p} x ^ {(n - p)} y \\ = \binom {m + n - p - 1 - h} {p} (y x ^ {(n - p)} + (n - p - 1 - h) x ^ {(n - p - 1)}) \\ = y \binom {m + n - p + 1 - h} {p} x ^ {(n - p)} \\ \qquad + \binom {m + n - p - 1 - h} {p} (n - p - 1 - h) x ^ {(n - p - 1)}. \end{array}
\]

Inserting this into (18), we obtain

\[
\begin{array}{l} (m + 1) (\operatorname{ad} x ^ {(n)}) y ^ {(m + 1)} \\ \qquad = \sum_ {p \geq 1} (m - p + 1) y ^ {(m - p + 1)} \binom {m + n - p + 1 - h} {p} x ^ {(n - p)} \\ \qquad + \sum_ {p \geq 1} y ^ {(m - p)} \binom {m + n - p - 1 - h} {p} (n - p - 1 - h) x ^ {(n - p - 1)} \\ \qquad + y ^ {(m)} (n - 1 - h) x ^ {(n - 1)} \\ \qquad = \sum_ {p \geq 1} (m - p + 1) y ^ {(m - p + 1)} \binom {m + n - p + 1 - h} {p} x ^ {(n - p)} \\ \qquad + \sum_ {p \geq 0} y ^ {(m - p)} \binom {m + n - p - 1 - h} {p} (n - p - 1 - h) x ^ {(n - p - 1)}. \end{array}
\]

Changing \(p\) to \(p - 1\) in the second sum, and regrouping the terms, we obtain

\[
(m + 1) (\operatorname{ad} x ^ {(n)}) y ^ {(m + 1)} = \sum_ {p \geq 1} y ^ {(m - p + 1)} A _ {p} x ^ {(n - p)}\tag{19}
\]

with

\[
A _ {p} = (m - p + 1) \binom{m + n - p + 1 - h}{p} + (n - p - h) \binom{m + n - p - h}{p - 1}.
\]

Putting \(z = m + n - p - h\) , this can also be written as

\[
\begin{array}{l} A _ {p} = \frac {1}{p !} (m - p + 1) (z + 1) z (z - 1) \dots (z - p + 2) \\ \qquad + \frac {1}{(p - 1) !} (z - m) z (z - 1) \dots (z - p + 2) \\ \qquad = \frac {1}{p !} z (z - 1) \dots (z - p + 2) [ (m - p + 1) (z + 1) + p (z - m) ] \\ \qquad = (m + 1) \binom {z} {p} = (m + 1) \binom {(m + 1) + n - p - 1 - h} {p}. \end{array}
\]

Inserting this into (19), we obtain \((17_{m + 1})\) , hence (i).

Assume that \(\binom{h}{p} \in \mathrm{A}'\) for \(p < n\) . Then, for all \(\mathrm{P} \in \mathbf{Q}[\mathrm{T}]\) of degree \(< n\) such that \(\mathrm{P}(\mathbf{Z}) \subset \mathbf{Z}\) , we have \(\mathrm{P}(h) \in \mathrm{A}'\) (no. 4, Cor. of Prop. 2). Hence, in view of (16) with \(m = n\) ,

\[
\begin{array}{c} (- 1) ^ {n} \binom{h}{n} = \binom{n - 1 - h}{n} \\ = - x ^ {(n)} y ^ {(n)} + \sum_ {p = 0} ^ {n - 1} y ^ {(n - p)} \binom{2 n - p - 1 - h}{p} x ^ {(n - p)} \in \mathrm{A} ^ {\prime}; \end{array}
\]

hence (ii) by induction on \(n\) .

\subsection{BIORDERS IN THE ENVELOPING ALGEBRA OF A SPLIT REDUCTIVE LIE ALGEBRA}

Let \(\mathfrak{g}\) be a reductive Lie algebra over \(\mathbf{Q}\) , \(\mathfrak{h}\) a splitting Cartan subalgebra of \(\mathfrak{g}\) , and \(R = R(\mathfrak{g},\mathfrak{h})\) (§2, no. 1, Remark 5).

DEFINITION 1. A lattice \(\mathcal{H}\) in \(\mathfrak{h}\) is said to be permissible (relative to \(\mathfrak{g}\) ) if, for all \(\alpha \in \mathrm{R}\) , we have \(H_{\alpha} \in \mathcal{H}\) and \(\alpha(\mathcal{H}) \subset \mathbf{Z}\) .

Remarks. 1) Let B be a basis of R. A lattice \(\mathcal{H}\) in \(\mathfrak{h}\) is permissible if and only if \(H_{\alpha} \in \mathcal{H}\) and \(\alpha(\mathcal{H}) \subset \mathbf{Z}\) for all \(\alpha \in B\) .

\begin{enumerate}
\def\labelenumi{\arabic{enumi})}
\setcounter{enumi}{1}
\item
  Let \(\mathfrak{c}\) be the centre of \(\mathfrak{g}\) . Then, a lattice \(\mathcal{H}\) in \(\mathfrak{h}\) is permissible if and only if \(\mathrm{Q}(\mathrm{R}^{\vee}) \subset \mathcal{H} \subset \mathrm{P}(\mathrm{R}^{\vee}) \oplus \mathfrak{c}\) . The lattice \(\mathcal{H} \cap \mathcal{D}\mathfrak{g}\) is then permissible in the Cartan subalgebra \(\mathfrak{h} \cap \mathcal{D}\mathfrak{g}\) of \(\mathcal{D}\mathfrak{g}\) . There may exist permissible lattices \(\mathcal{H}\) such that \(\mathcal{H} \neq (\mathcal{H} \cap \mathcal{D}\mathfrak{g}) \oplus (\mathcal{H} \cap \mathfrak{c})\) (cf.~§13, no. 1.IX).
\item
  If \(\mathfrak{g}\) is semi-simple, the permissible lattices in \(\mathfrak{h}\) are the subgroups \(\mathcal{H}\) of \(\mathfrak{h}\) such that \(Q(R^{\vee})\subset \mathcal{H}\subset P(R^{\vee})\) .
\end{enumerate}

In the remainder of this number, we assume fixed a split reductive Lie algebra \((\mathfrak{g},\mathfrak{h})\) , a basis B of \(\mathrm{R}=\mathrm{R}(\mathfrak{g},\mathfrak{h})\) and, for each \(\alpha\in B\) , a pair \((x_{\alpha},y_{\alpha})\) with

\[
y _ {\alpha} \in \mathfrak {g} ^ {- \alpha}, x _ {\alpha} \in \mathfrak {g} ^ {\alpha}, [ y _ {\alpha}, x _ {\alpha} ] = H _ {\alpha}.\tag{20}
\]

If we denote by \(n_{+}\) (resp. \(n_{-}\) ) the subalgebra of g generated by the \(x_{\alpha}\) (resp. the \(y_{\alpha}\) ), we know (§3, no. 3, Prop. 9 (iii)) that

\[
\mathfrak {g} = \mathfrak {n} _ {-} \oplus \mathfrak {h} \oplus \mathfrak {n} _ {+}\tag{21}
\]

\[
\mathrm{U} (\mathfrak {g}) = \mathrm{U} (\mathfrak {n} _ {-}) \otimes_ {\mathbf {Q}} \mathrm{U} (\mathfrak {h}) \otimes_ {\mathbf {Q}} \mathrm{U} (\mathfrak {n} _ {+})\tag{22}
\]

(where \(\mathrm{U}(\mathfrak{g}),\ldots\) are the enveloping algebras of \(\mathfrak{g},\ldots\) ).

Denote by \(\mathcal{U}_{+}\) the \(\mathbf{Z}\)-subalgebra of \(\mathrm{U}(\mathfrak{n}_{+})\) generated by the \(x_{\alpha}^{(n)}\) for \(\alpha \in B\) and \(n \in \mathbf{N}\) . Let W be the Weyl group of R, \(R_{+}\) the set of positive roots relative to B.

Lemma 5. (i) \(\mathcal{U}_{+}\) is a lattice in the vector space \(\mathrm{U}(\mathfrak{n}_{+})\) .

\begin{enumerate}
\def\labelenumi{(\roman{enumi})}
\setcounter{enumi}{1}
\tightlist
\item
  For all \(\alpha \in \mathrm{B}\) , we have \(\mathcal{U}_{+} \cap \mathrm{U}(\mathfrak{g}^{\alpha}) = \bigoplus_{n \in \mathbf{N}} \mathbf{Z}x_{\alpha}^{(n)}\) .
\end{enumerate}

By definition, \(\mathcal{U}_{+}\) is generated as a \(\mathbf{Z}\) -module by the elements

\[
x _ {\varphi} ^ {(\mathbf {n})} = \prod_ {1 \leq i \leq r} x _ {\varphi (i)} ^ {(n (i))}
\]

where \(r \in \mathbf{N}\) , \(\varphi = (\varphi(i)) \in \mathrm{B}^{r}\) , and \(\mathbf{n} = (n(i)) \in \mathbf{N}^{r}\) . Give the algebra \(\mathrm{U}(\mathfrak{n}_{+})\) the graduation of type \(\mathrm{Q}(\mathrm{R})\) for which each \(\mathfrak{g}^{\alpha} (\alpha \in \mathrm{R}_{+})\) is homogeneous of degree \(\alpha\) . A monomial \(x_{\varphi}^{(\mathbf{n})}\) of the preceding type is homogeneous of degree

\[
\sum_ {1 \leq i \leq r} n (i) \varphi (i) \in \mathrm{Q} (\mathrm{R}).
\]

The monomials of this kind having a given degree q are finite in number, and generate over Q the homogeneous component of \(\mathrm{U}(\mathfrak{n}_{+})\) of degree q. This proves (i).

If \(\alpha \in \mathrm{B}\) , \(\mathcal{U}_{+} \cap \mathrm{U}(\mathfrak{g}^{\alpha})\) is contained in the sum of the homogeneous components of degrees which are multiples of \(\alpha\) ; thus, by the preceding, \(\mathcal{U}_{+} \cap \mathrm{U}(\mathfrak{g}^{\alpha})\) is generated by the \(x_{\varphi}^{(\mathbf{n})}\) such that \(\sum n(i)\varphi(i) \in \mathbf{N}\alpha\) , which forces \(\varphi(i) = \alpha\) for all \(i\) (since B is a basis of R), so

\[
x _ {\varphi} ^ {(\mathbf {n})} = x _ {\alpha} ^ {(n (1))} \dots x _ {\alpha} ^ {(n (r))} = \frac {(n (1) + \cdots + n (r)) !}{n (1) ! \dots n (r) !} x _ {\alpha} ^ {(n (1) + \dots + n (r))}.
\]

Thus, \(\mathcal{U}_{+} \cap \mathcal{U}(\mathfrak{g}^{\alpha}) \subset \bigoplus_{n} \mathbf{Z}x_{\alpha}^{(n)}\) , hence (ii).

Q.E.D.

In the remainder of this paragraph, if E and F are Z-submodules of \(U(\mathfrak{g})\) , we denote by E.F the Z-submodule of \(U(\mathfrak{g})\) generated by the products ab, where \(a \in E, b \in F\) .

PROPOSITION 3. Let \(\mathcal{H}\) be a permissible lattice in \(\mathfrak{h}\) . Let \(\mathcal{U}_{+},\mathcal{U}_{-},\mathcal{U}_{0}\) be the \(\mathbf{Z}\) -subalgebras of \(\mathrm{U}(\mathfrak{g})\) generated respectively by the elements \(x_{\alpha}^{(n)}\) ( \(\alpha \in \mathrm{B}, n \in \mathbf{N}\) ), \(y_{\alpha}^{(n)}\) ( \(\alpha \in \mathrm{B}, n \in \mathbf{N}\) ), \(\binom{h}{n}\) ( \(h \in \mathcal{H}, n \in \mathbf{N}\) ). Let \(\mathcal{U}\) be the \(\mathbf{Z}\) -subalgebra of \(\mathrm{U}(\mathfrak{g})\) generated by \(\mathcal{U}_{+},\mathcal{U}_{-},\mathcal{U}_{0}\) .

\begin{enumerate}
\def\labelenumi{(\roman{enumi})}
\item
  \(\mathcal{U}\) is a biorder in the bigebra \(\mathrm{U}(\mathfrak{g})\) .
\item
  We have \(\mathcal{U} = \mathcal{U}_{-}. \mathcal{U}_{0}. \mathcal{U}_{+}\) , \(\mathcal{U} \cap \mathfrak{h} = \mathcal{H}\) and, for all \(\alpha \in B\) ,
\end{enumerate}

\[
\mathcal {U} \cap \mathfrak {g} ^ {\alpha} = \mathbf {Z} x _ {\alpha}, \quad \mathcal {U} \cap \mathfrak {g} ^ {- \alpha} = \mathbf {Z} y _ {\alpha}.
\]

By Lemma 5 and Prop. 2, \(\mathcal{U}_{+},\mathcal{U}_{-},\mathcal{U}_{0}\) are orders in the \(\mathbf{Q}\) -algebras \(\mathrm{U}(\mathfrak{n}_{+}),\mathrm{U}(\mathfrak{n}_{-}),\mathrm{U}(\mathfrak{h})\) , respectively, and

\[
\binom{\pm h + q}{p} \in \mathcal {U} _ {0} \quad \text { for } h \in \mathcal {H}, q \in \mathbf {Z}, p \in \mathbf {N}.\tag{23}
\]

Put \(\mathcal{L} = \mathcal{U}_{-}. \mathcal{U}_{0}. \mathcal{U}_{+} \subset \mathrm{U}(\mathfrak{g})\) . By (22), \(\mathcal{L}\) is a lattice in \(\mathrm{U}(\mathfrak{g})\) . By construction,

\[
\mathcal {U} _ {-}. \mathcal {L} \subset \mathcal {L}\tag{24}
\]

\[
\mathcal {L}. \mathcal {U} _ {+} \subset \mathcal {L}\tag{25}
\]

while Lemma 3 and (23) imply that

\[
\mathcal {U} _ {0}. \mathcal {L} \subset \mathcal {L}\tag{26}
\]

\[
\mathcal {L}. \mathcal {U} _ {0} \subset \mathcal {L}.\tag{27}
\]

\[
\text {   Let   } \alpha \in \mathrm{B}, n \in \mathbf {N}, r \in \mathbf {N}, \varphi = (\varphi (i)) \in \mathrm{B} ^ {r}, \text {   and   }
\]

\[
(m (1), \dots , m (r)) \in \mathbf {N} ^ {r}.
\]

We show that

\[
x _ {\alpha} ^ {(n)} y _ {\varphi (1)} ^ {(m (1))} \dots y _ {\varphi (r)} ^ {(m (r))} \in \mathcal {L}\tag{28}
\]

or equivalently, in view of (25), that

\[
[ x _ {\alpha} ^ {(n)}, y _ {\varphi (1)} ^ {(m (1))} \dots y _ {\varphi (r)} ^ {(m (r))} ] \in \mathcal {L}.\tag{29}
\]

We argue by induction on \(r\) . The bracket to be studied is the sum of the terms

\[
y _ {\varphi (1)} ^ {(m (1))} \dots y _ {\varphi (k)} ^ {(m (k))} [ x _ {\alpha} ^ {(n)}, y _ {\varphi (k + 1)} ^ {(m (k + 1))} ] y _ {\varphi (k + 2)} ^ {(m (k + 2))} \dots y _ {\varphi (r)} ^ {(m (r))}.\tag{30}
\]

For \(\alpha \neq \varphi (k + 1)\) , \(x_{\alpha}\) and \(y_{\varphi (k + 1)}\) commute, so \([x_{\alpha}^{(n)},y_{\varphi (k + 1)}^{(m(k + 1))}] = 0\) . If \(\alpha = \varphi (k + 1)\) , the expression (30) is, by (17), the sum of expressions of the form

\[
y _ {\varphi (1)} ^ {(m (1))} \dots y _ {\varphi (k)} ^ {(m (k))} y _ {\varphi (k + 1)} ^ {(m (k + 1) - p)} \binom{q - h}{p} x _ {\alpha} ^ {(n - p)} y _ {\varphi (k + 2)} ^ {(m (k + 2))} \dots y _ {\varphi (r)} ^ {(m (r))}\tag{31}
\]

where \(q \in \mathbf{Z}, p \in \mathbf{N} - \{0\}, h \in \mathcal{H}\) . The induction hypothesis, together with (24) and (26), proves that the expression (31) belongs to \(\mathcal{L}\) . We have thus proved (28).

By (28), \(x_{\alpha}^{(n)}\mathcal{U}_{-}\subset \mathcal{L}\) ; thus, by (25) and (27), \(x_{\alpha}^{(n)}\mathcal{L}\subset \mathcal{L}\) , so \(\mathcal{U}_{+}.\mathcal{L}\subset \mathcal{L}\) and

\[
\mathcal {L}. \mathcal {L} \subset \mathcal {U} _ {-}. \mathcal {U} _ {0}. \mathcal {L} \subset \mathcal {U} _ {-}. \mathcal {L} \subset \mathcal {L}.
\]

Thus, \(\mathcal{L}\) is a \(\mathbf{Z}\) -subalgebra of \(\mathrm{U}(\mathfrak{g})\) , so \(\mathcal{U} = \mathcal{L}\) . If \(c\) is the coproduct of \(\mathrm{U}(\mathfrak{g})\) , \(c(\mathcal{U}) \subset \mathcal{U} \otimes_{\mathbf{Z}} \mathcal{U}\) (no. 2, Prop. 1). Let \(\gamma\) be the counit of \(\mathrm{U}(\mathfrak{g})\) . Since \(\gamma(x_{\alpha}^{(n)}) = \gamma(y_{\alpha}^{(n)}) = \gamma\left(\binom{h}{n}\right) = 0\) for \(n > 0\) , we have \(\gamma(\mathcal{U}) \subset \mathbf{Z}\) . This proves (i). On the other hand,

\[
\mathcal {U} \cap \mathfrak {h} = \mathcal {L} \cap \mathfrak {h} = \mathcal {U} _ {0} \cap \mathfrak {h} = \mathcal {H}
\]

by Prop. 2 of no. 4; similarly,

\[
\mathcal {U} \cap \mathfrak {g} ^ {\alpha} = \mathcal {U} _ {+} \cap \mathfrak {g} ^ {\alpha} = \mathbf {Z} x _ {\alpha}
\]

by Lemma 5. This proves (ii).

Remark 4. By Prop. 5 of §4, no. 4, there exists a unique automorphism \(\theta\) of \(\mathfrak{g}\) such that \(\theta(x_{\alpha}) = y_{\alpha}\) and \(\theta(y_{\alpha}) = x_{\alpha}\) for all \(\alpha \in \mathrm{B}\) , and \(\theta(h) = -h\) for all \(h \in \mathfrak{h}\) ; we have \(\theta^2 = 1\) . By construction of \(\mathcal{U}\) , we see that the automorphism of \(\mathrm{U}(\mathfrak{g})\) that extends \(\theta\) leaves \(\mathcal{U}\) stable.

COROLLARY 1. Put \(\mathcal{G} = \mathcal{U} \cap \mathfrak{g}\) . Then \(\mathcal{G}\) is an order in the Lie algebra \(\mathfrak{g}\) , stable under \(\theta\) . We have \(\mathcal{G} = \mathcal{H} + \sum_{\alpha \in \mathrm{R}} (\mathcal{G} \cap \mathfrak{g}^{\alpha})\) . For all \(\alpha \in \mathrm{B}\) and all \(n \in \mathbf{N}\) , the maps \((\operatorname{ad} x_{\alpha})^{n} / n!\) , \((\operatorname{ad} y_{\alpha})^{n} / n!\) leave \(\mathcal{U}\) and \(\mathcal{G}\) stable.

The first assertion is clear. The second follows by considering the graduation of type \(Q(R)\) on \(U(\mathfrak{g})\) and U. The third follows from Lemma 2 of no. 5.

COROLLARY 2. Let \(w \in \mathrm{W}\) . There exists an elementary automorphism \(\varphi\) of \(\mathfrak{g}\) that commutes with \(\theta\) , leaves \(\mathcal{G}\) and \(\mathcal{U}\) stable, and extends \(w\) .

It suffices to treat the case in which \(w\) is of the form \(s_{\alpha}\) ( \(\alpha \in \mathrm{B}\) ). Note first of all that \(\operatorname{ad} x_{\alpha}\) and \(\operatorname{ad} y_{\alpha}\) are locally nilpotent on \(\mathrm{U}(\mathfrak{g})\) , in other words that for all \(u \in \mathrm{U}(\mathfrak{g})\) there exists an integer \(n\) such that \((\operatorname{ad} x_{\alpha})^n u = (\operatorname{ad} y_{\alpha})^n u = 0\) . This enables us to define the automorphisms \(e^{\operatorname{ad} x_{\alpha}} = \sum_{n=0}^{\infty} \frac{1}{n!} (\operatorname{ad} x_{\alpha})^n\) and \(e^{\mathrm{ad}y_{\alpha}}\) of \(\mathrm{U}(\mathfrak{g})\) ; we verify immediately that these automorphisms of \(\mathrm{U}(\mathfrak{g})\) leave \(\mathcal{U}\) stable. Put \(\varphi_1 = e^{\mathrm{ad}x_\alpha}e^{\mathrm{ad}y_\alpha}e^{\mathrm{ad}x_\alpha}\) , \(\varphi_2 = e^{\mathrm{ad}y_\alpha}e^{\mathrm{ad}x_\alpha}e^{\mathrm{ad}y_\alpha}\) . We have \(\varphi_1|\mathfrak{g} = \varphi_2|\mathfrak{g}\) (§2, no. 2, formula (1)), so \(\varphi_1 = \varphi_2\) . Put \(\varphi_1 = \varphi_2 = \varphi\) . We have \(\theta \varphi \theta^{-1} = \varphi\) , so \(\theta\) and \(\varphi\) commute. On the other hand, \(\varphi |\mathfrak{h} = w\) by §2, no. 2, Lemma 1.

COROLLARY 3. Let \(\alpha \in \mathrm{R}\) . If \(x \in \mathcal{G} \cap \mathfrak{g}^{\alpha}\) and \(n \in \mathbf{N}\) , we have \(x^{(n)} \in \mathcal{U}\) , and \((\operatorname{ad} x)^n / n!\) leaves \(\mathcal{G}\) and \(\mathcal{U}\) stable.

This is clear if \(\alpha \in \mathrm{B}\) , by construction of \(\mathcal{U}\) and Cor. 1. In the general case, there exists \(w \in \mathrm{W}\) such that \(w(\alpha) \in \mathrm{B}\) (Chap. VI, §1, no. 5, Prop. 15). By Cor. 2, there exists an automorphism \(\varphi\) of \(\mathfrak{g}\) that leaves \(\mathcal{G}\) and \(\mathcal{U}\) stable and takes \(\mathfrak{g}^{\alpha}\) to \(\mathfrak{g}^{w(\alpha)}\) , hence the corollary by transport of structure.

COROLLARY 4. There exists a Chevalley system \((X_{\alpha})_{\alpha \in \mathrm{R}}\) in \((\mathfrak{g},\mathfrak{h})\) (§2, no. 4, Def. 3) such that \(X_{\alpha} = x_{\alpha}\) and \(X_{-\alpha} = y_{\alpha}\) for \(\alpha \in \mathrm{B}\) . For every Chevalley system \((X_{\alpha}^{\prime})_{\alpha \in \mathrm{R}}\) having these properties, and for all \(\alpha \in \mathrm{R}\) , \(X_{\alpha}^{\prime}\) is a basis of \(\mathcal{G} \cap \mathfrak{g}^{\alpha}\) .

For \(\alpha \in \mathrm{B}\) , put \(X_{\alpha} = x_{\alpha}, X_{-\alpha} = y_{\alpha}\) . For \(\alpha \in \mathrm{R}_{+} - \mathrm{B}\) , choose a \(w \in \mathrm{W}\) such that \(w(\alpha) \in \mathrm{B}\) and an automorphism \(\varphi\) of \(\mathfrak{g}\) such that \(\theta \varphi = \varphi \theta\) , \(\varphi(\mathcal{G}) = \mathcal{G}\) and \(\varphi(h) = w^{-1}(h)\) for \(h \in \mathfrak{h}\) (Cor. 2); put \(X_{\alpha} = \varphi(x_{w(\alpha)})\) , \(X_{-\alpha} = \varphi(y_{w(\alpha)})\) . Then

\[
[ X _ {- \alpha}, X _ {\alpha} ] = \varphi ([ y _ {w (\alpha)}, x _ {w (\alpha)} ]) = \varphi (H _ {w (\alpha)}) = w ^ {- 1} (H _ {w (\alpha)}) = H _ {\alpha}
\]

\[
\theta (X _ {\alpha}) = \theta \varphi (x _ {w (\alpha)}) = \varphi \theta (x _ {w (\alpha)}) = \varphi (y _ {w (\alpha)}) = X _ {- \alpha}
\]

so \((X_{\alpha})_{\alpha \in \mathrm{R}}\) is a Chevalley system. Moreover,

\[
\mathcal {G} \cap \mathfrak {g} ^ {\alpha} = \varphi (\mathcal {G} \cap \mathfrak {g} ^ {w (\alpha)}) = \varphi (\mathbf {Z} x _ {w (\alpha)}) = \mathbf {Z} X _ {\alpha}\tag{32}
\]

\[
\mathcal {G} \cap \mathfrak {g} ^ {- \alpha} = \varphi (\mathcal {G} \cap \mathfrak {g} ^ {- w (\alpha)}) = \varphi (\mathbf {Z} y _ {w (\alpha)}) = \mathbf {Z} X _ {- \alpha}.\tag{33}
\]

Let \((X'_{\alpha})_{\alpha \in \mathrm{R}}\) be a Chevalley system such that \(X'_{\alpha} = x_{\alpha}, X'_{-\alpha} = y_{\alpha}\) for \(\alpha \in \mathrm{B}\) . Let \(\mathrm{S}\) be the set of \(\alpha \in \mathrm{R}\) such that \(X'_{\alpha} = \pm X_{\alpha}\) . By §2, no. 4, Prop. 7, \(\mathrm{S}\) is a closed set of roots. Since \(\mathrm{S} \supset \mathrm{B} \cup (-\mathrm{B})\) , we have \(\mathrm{S} = \mathrm{R}\) (Chap. VI, §1, no. 6, Prop. 19). Thus, by (32) and (33), we have \(\mathcal{G} \cap \mathfrak{g}^{\alpha} = \mathbf{Z}X'_{\alpha}\) for all \(\alpha \in \mathrm{R}\) .

Remarks. 5) Let \((X_{\alpha})_{\alpha \in \mathrm{R}}\) be the Chevalley system constructed above. If \(\alpha, \beta, \alpha + \beta \in \mathrm{R}\) and if we put \([X_{\alpha}, X_{\beta}] = \mathrm{N}_{\alpha, \beta} X_{\alpha + \beta}\) , we have \([X_{\alpha}, X_{\beta}] \in \mathcal{G} \cap \mathfrak{g}^{\alpha + \beta}\) , and we recover the fact that \(\mathrm{N}_{\alpha, \beta} \in \mathbf{Z}\) (cf.~§2, no. 4, Prop. 7).

\begin{enumerate}
\def\labelenumi{\arabic{enumi})}
\setcounter{enumi}{5}
\tightlist
\item
  We have obtained in passing a new proof of the existence of Chevalley systems (cf.~§4, no. 4, Cor. of Prop. 5), independent of Lemma 4, §2.
\end{enumerate}

\subsection{CHEVALLEY ORDERS}

Let \((\mathfrak{g},\mathfrak{h})\) be a split reductive Lie algebra over \(\mathbf{Q}\) , \(\mathrm{R}\) its root system. Choose:

\begin{enumerate}
\def\labelenumi{\alph{enumi})}
\item
  a permissible lattice \(\mathcal{H}\) in \(\mathfrak{h}\) (no. 6, Def. 1);
\item
  for all \(\alpha \in \mathrm{R}\) , a lattice \(\mathcal{G}^{\alpha}\) in \(\mathfrak{g}^{\alpha}\) .
\end{enumerate}

Put \(\mathcal{G} = \mathcal{H} \oplus \sum_{\alpha \in \mathrm{R}} \mathcal{G}^{\alpha}\) . This is a lattice in \(\mathfrak{g}\) . Denote by \(\mathcal{U}\) the \(\mathbf{Z}\) -subalgebra of \(\mathrm{U}(\mathfrak{g})\) generated by the \(\binom{h}{n}\) ( \(h \in \mathcal{H}, n \in \mathbf{N}\) ) and the \(x^{(n)}\) ( \(x \in \mathcal{G}^{\alpha}, \alpha \in \mathrm{R}, n \in \mathbf{N}\) ). Finally, for \(\alpha \in \mathrm{R}\) and \(x \in \mathfrak{g}^{\alpha} - \{0\}\) , put

\[
w _ {\alpha} (x) = (\exp \operatorname{ad} x) (\exp \operatorname{ad} y) (\exp \operatorname{ad} x),
\]

where \(y\) is the unique element of \(\mathfrak{g}^{-\alpha}\) such that \([y,x] = H_{\alpha}\) . With these notations:

THEOREM 2. The following conditions are equivalent:

\begin{enumerate}
\def\labelenumi{(\roman{enumi})}
\item
  There exists a Chevalley system \((X_{\alpha})_{\alpha \in \mathrm{R}}\) of \((\mathfrak{g},\mathfrak{h})\) such that \(\mathcal{G}_{\alpha} = \mathbf{Z}X_{\alpha}\) for all \(\alpha \in \mathrm{R}\) .
\item
  \(\mathcal{U} \cap \mathfrak{g} = \mathcal{G}\) and \([\mathcal{G}^{\alpha}, \mathcal{G}^{-\alpha}] = \mathbf{Z}H_{\alpha}\) for all \(\alpha \in \mathrm{R}\) .
\item
  For all \(\alpha \in \mathrm{R}, x \in \mathcal{G}^{\alpha}, n \in \mathbf{N}\) , the endomorphism \((\operatorname{ad} x)^n / n!\) of \(\mathfrak{g}\) maps \(\mathcal{G}\) to \(\mathcal{G}\) , and \([\mathcal{G}^{\alpha}, \mathcal{G}^{-\alpha}] = \mathbf{Z}H_{\alpha}\) .
\item
  For all \(\alpha \in \mathrm{R}\) and every basis \(x\) of \(\mathcal{G}^{\alpha}\) , \(w_{\alpha}(x)\) maps \(\mathcal{G}\) to \(\mathcal{G}\) (that is, maps \(\mathcal{G}^{\beta}\) to \(\mathcal{G}^{s_{\alpha}(\beta)}\) for all \(\beta \in \mathrm{R}\) ).
\item
  \(\Longrightarrow\) (ii): let \((X_{\alpha})_{\alpha \in \mathrm{R}}\) be a Chevalley system in \((\mathfrak{g},\mathfrak{h})\) such that \(\mathcal{G}^{\alpha} = \mathbf{Z}X_{\alpha}\) for all \(\alpha \in \mathrm{R}\) , and let B be a basis of R. For \(\alpha \in \mathrm{B}\) , put \(x_{\alpha} = X_{\alpha},y_{\alpha} = X_{-\alpha}\) . Let \(\mathcal{U}'\) be the biorder associated by Prop. 3 of no. 6 to \(\mathcal{H}\) , the \(x_{\alpha}\) and the \(y_{\alpha}\) . It is clear that \(\mathcal{U}' \subset \mathcal{U}\) . By Cor. 3 and 4 of Prop. 3, \(x^{(n)} \in \mathcal{U}'\) for all \(\alpha \in \mathrm{R}\) , \(x \in \mathcal{G}^{\alpha}\) and \(n \in \mathbf{N}\) . Thus \(\mathcal{U} = \mathcal{U}'\) , which proves (ii).
\item
  \(\Longrightarrow\) (iii): this is clear by Lemma 2 of no. 5.
\item
  \(\Longrightarrow\) (iv): let \(\alpha \in \mathrm{R}\) and let \(x\) be a basis of \(\mathcal{G}^{\alpha}\) . Since \([\mathcal{G}^{\alpha},\mathcal{G}^{-\alpha}] = \mathbf{Z}H_{\alpha}\) , the unique \(y \in \mathfrak{g}^{-\alpha}\) such that \([y,x] = H_{\alpha}\) belongs to \(\mathcal{G}^{-\alpha}\) . Since \(\exp \mathrm{ad} x\) and \(\exp \mathrm{ad} y\) leave \(\mathcal{G}\) stable by (iii), so does \(w_{\alpha}(x)\) .
\item
  \(\Longrightarrow\) (i): let B be a basis of R. Choose a basis \(x_{\alpha}\) of \(\mathcal{G}^{\alpha}\) for all \(\alpha \in \mathrm{B}\) . Let \(y_{\alpha} \in \mathcal{G}^{-\alpha}\) be such that \([y_{\alpha}, x_{\alpha}] = H_{\alpha}\) . By §1, no. 5, formulas (5), we have \(y_{\alpha} = w_{\alpha}(x_{\alpha}). x_{\alpha}\) so \(y_{\alpha}\) is a basis of \(\mathcal{G}^{-\alpha}\) by (iv). Let \(\mathcal{G}\) be the order in \(\mathfrak{g}\) defined by \(\mathcal{H}\) , the \(x_{\alpha}\) and the \(y_{\alpha}\) (no. 6, Cor. 1 of Prop. 3). Then \(\mathcal{G}\) is stable under the \((\operatorname{ad} x_{\alpha})^{n} / n!\) , \((\operatorname{ad} y_{\alpha})^{n} / n!\) (loc. cit.), and hence under the \(w_{\alpha}(x_{\alpha})\) .
\end{enumerate}

Now let \(\beta \in \mathrm{R}\) . There exist \(\alpha_0, \alpha_1, \ldots, \alpha_r \in \mathrm{B}\) such that

\[
\beta = s _ {\alpha_ {r}} s _ {\alpha_ {r - 1}} \dots s _ {\alpha_ {1}} (\alpha_ {0})
\]

(Chap. VI, §1, no. 5, Prop. 15). Then \(w_{\alpha_r}(x_{\alpha_r}).w_{\alpha_{r-1}}(x_{\alpha_{r-1}}) \ldots w_{\alpha_1}(x_{\alpha_1})\) maps \(\mathcal{G}^{\alpha_0}\) to \(\mathcal{G}^\beta\) by (iv), and maps \(\mathcal{G} \cap \mathfrak{g}^{\alpha_0}\) to \(\mathcal{G} \cap \mathfrak{g}^\beta\) by the preceding. Since \(\mathcal{G} \cap \mathfrak{g}^{\alpha_0} = \mathcal{G}^{\alpha_0}\) (Prop. 3 (ii)), we have \(\mathcal{G} \cap \mathfrak{g}^\beta = \mathcal{G}^\beta\) . Thus

\[
\mathcal {G} = \mathcal {H} \oplus \sum_ {\beta \in \mathrm{R}} (\mathcal {G} \cap \mathfrak {g} ^ {\beta}) = \mathcal {H} \oplus \sum_ {\beta \in \mathrm{R}} \mathcal {G} ^ {\beta} = \mathcal {G}
\]

and Cor. 4 of Prop. 3 concludes the proof.

DEFINITION 2. When conditions (i) to (iv) of Th. 2 are satisfied, \(\mathcal{G}\) is said to be a Chevalley order in \((\mathfrak{g}, \mathfrak{h})\) .

Remark. Chevalley orders in \((\mathfrak{g},\mathfrak{h})\) always exist. Indeed, the Chevalley orders are the sets of the form \(\mathcal{H}\oplus\sum_{\alpha\in \mathrm{R}}\mathbf{Z}X_{\alpha}\) , where \((X_{\alpha})_{\alpha\in \mathrm{R}}\) is a Chevalley system in \((\mathfrak{g},\mathfrak{h})\) and \(\mathcal{H}\) is a lattice in \(\mathfrak{h}\) such that

\[
\mathrm{Q} (\mathrm{R} ^ {\vee}) \subset \mathcal {H} \subset \mathrm{P} (\mathrm{R} ^ {\vee}) \oplus \mathfrak {c}
\]

(c being the centre of g).

THEOREM 3. We retain the notations at the beginning of no. 7, and assume that \(\mathcal{G}\) is a Chevalley order in \((\mathfrak{g},\mathfrak{h})\) .

\begin{enumerate}
\def\labelenumi{(\roman{enumi})}
\item
  \(\mathcal{U}\) is a b-order in \(\mathrm{U}(\mathfrak{g})\) .
\item
  Let B be a basis of \(\mathrm{R}\), and \((X_{\alpha})_{\alpha\in\mathrm{B}\cup(-\mathrm{B})}\) a family of elements of \(\mathfrak{g}\) such that \(\mathcal{G}^{\alpha}=\mathbf{Z}X_{\alpha}\) for \(\alpha\in\mathrm{B}\cup(-\mathrm{B})\) . The \(\mathbf{Z}\)-algebra \(\mathcal{U}\) is generated by the \(\binom{h}{n}\) and the \(X_{\alpha}^{(n)}\) ( \(h\in\mathcal{H},\alpha\in\mathrm{B}\cup(-\mathrm{B}),n\in\mathbf{N}\) ). If \(\mathfrak{g}\) is semi-simple and \(\mathcal{H}=\mathrm{Q}(\mathrm{R}^{\vee})\) , the \(\mathbf{Z}\)-algebra \(\mathcal{U}\) is generated by the \(X_{\alpha}^{(n)}\) ( \(\alpha\in\mathrm{B}\cup(-\mathrm{B}),n\in\mathbf{N}\) ).
\item
  Let B be a basis of R, \(R_{+}\) the corresponding set of positive roots, \(R_{-} = -R_{+}\) , \(n_{+} = \sum_{\alpha \in R_{+}} g^{\alpha}\) , \(n_{-} = \sum_{\alpha \in R_{-}} g^{\alpha}\) . Then,
\end{enumerate}

\[
\mathcal {U} = (\mathcal {U} \cap \mathrm{U} (\mathfrak {n} _ {-})). (\mathcal {U} \cap \mathrm{U} (\mathfrak {h})). (\mathcal {U} \cap \mathrm{U} (\mathfrak {n} _ {+})).
\]

Let \((h_i)_{i\in \mathrm{I}}\) be a basis of \(\mathcal{H}\) . For all \(\alpha \in \mathrm{R}\) , let \(X_{\alpha}\) be a basis of \(\mathcal{G}^{\alpha}\) . Give the set \(\mathrm{I}\cup \mathrm{R}\) a total order (we assume that \(\mathrm{I}\cap \mathrm{R} = \emptyset\) ). For \(\lambda \in \mathrm{I}\cup \mathrm{R}\) and \(n\in \mathbf{N}\) , put \(e_{\lambda}^{\langle n\rangle} = \binom{h_{\lambda}}{n}\) if \(\lambda \in \mathrm{I}\) , \(e_{\lambda}^{\langle n\rangle} = X_{\lambda}^{(n)}\) if \(\lambda \in \mathrm{R}\) . Then the products \(\prod_{\lambda \in \mathrm{I}\cup \mathrm{R}}e_{\lambda}^{\langle n_{\lambda}\rangle}\) , where \((n_{\lambda})\) belongs to \(\mathbf{N}^{\mathrm{I}\cup \mathrm{R}}\) , form a basis of the \(\mathbf{Z}\) -module \(\mathcal{U}\) . The products \(\prod_{\lambda \in \mathrm{I}}\binom{h_{\lambda}}{n_{\lambda}}\) , where \((n_{\lambda})\) belongs to \(\mathbf{N}^{\mathrm{I}}\) , form a basis of the \(\mathbf{Z}\) -module \(\mathcal{U}\cap \mathrm{U}(\mathfrak{h})\) . The products \(\prod_{\lambda \in \mathrm{R}_{+}}X_{\lambda}^{(n_{\lambda})}\) , where \((n_{\lambda})\) belongs to \(\mathbf{N}^{\mathrm{R}_{+}}\) , form a basis of the \(\mathbf{Z}\) -module \(\mathcal{U}\cap \mathrm{U}(\mathfrak{n}_{+})\) .

Let B and \((X_{\alpha})_{\alpha\in\mathrm{B}\cup(-\mathrm{B})}\) be as in (ii), and such that \([X_{-\alpha},X_{\alpha}]=H_{\alpha}\) . Let \(\mathcal{U}'\) be the \(\mathbf{Z}\)-subalgebra of \(\mathrm{U}(\mathfrak{g})\) generated by the \(\binom{h}{n}\) and the \(X_{\alpha}^{(n)}\) \((h\in\mathcal{H},\alpha\in\mathrm{B}\cup(-\mathrm{B}),n\in\mathbf{N})\) . We have seen in the proof of Th. 2, (i) \(\Longrightarrow\) (ii), that \(\mathcal{U}'\) is equal to \(\mathcal{U}\) and is a biorder in \(\mathrm{U}(\mathfrak{g})\) . This proves (i) and the first assertion of (ii); the second follows from Lemma 4 (ii). Assertion (iii) follows from Th. 1 (no. 3) and Prop. 3 (no. 6).

\subsection{ADMISSIBLE LATTICES}

Generalizing the terminology adopted for vector spaces, an endomorphism u of a module M is said to be diagonalizable if there exists a basis of M such that the matrix of u relative to this basis is diagonal.

Lemma 6. Let M be a free \(\mathbf{Z}\)-module of finite type, u an endomorphism of M, and v the endomorphism \(u \otimes 1\) of \(M \otimes_{\mathbf{Z}} \mathbf{Q}\) . Assume that \(\binom{v}{n}(M) \subset M\) for all \(n \in \mathbf{N}\) . Then u is diagonalizable.

\begin{enumerate}
\def\labelenumi{\alph{enumi})}
\item
  For any polynomial \(\mathrm{P} \in \mathbf{Q}[\mathrm{T}]\) such that \(\mathrm{P}(\mathbf{Z}) \subset \mathbf{Z}\) , we have \(\mathrm{P}(v)(\mathrm{M}) \subset \mathrm{M}\) (no. 4, Cor. of Prop. 2), so det \(\mathrm{P}(v) \in \mathbf{Z}\) .
\item
  Denote by \(\chi_{v}(t)=t^{d}+\alpha_{1}t^{d-1}+\cdots\) the characteristic polynomial of v. Let \(k\in\mathbf{Z}, n\in\mathbf{N}\) . Applying a) to the polynomial \(\binom{T-k}{n}\) , we see that the number
\end{enumerate}

\[
\begin{array}{l} a _ {n} = \det \binom {v - k} {n} = \frac {1}{(n !) ^ {d}} \det (v - k) \det (v - k - 1) \dots \det (v - k - n + 1) \\ = \frac {(- 1) ^ {n}}{(n !) ^ {4}} \chi_ {v} (k) \chi_ {v} (k + 1) \dots \chi_ {v} (k + n - 1) \end{array}
\]

is an integer. Take \(k - 1 < -\alpha_{1} / d\) . Then

\[
\chi_ {v} (k + n - 1) = n ^ {d} + (\alpha_ {1} + (k - 1) d) n ^ {d - 1} + \dots
\]

and

\[
| a _ {n} | = \frac {| \chi_ {v} (k + n - 1) |}{n ^ {d}} | a _ {n - 1} |;
\]

hence, if \(a_{n} \neq 0\) for all \(n \in \mathbf{N}\) , the sequence of the \(|a_{n}|\) is strictly decreasing for n sufficiently large, which is absurd. It follows that v has an integer eigenvalue \(\lambda\) . Put \(\mathrm{M}' = \mathrm{Ker}(u - \lambda.1)\) and \(M'' = M/M'\) . Then \(M'\) is the intersection with M of a vector subspace of \(M \otimes_{\mathbf{Z}} \mathbf{Q}\) , so the \(\mathbf{Z}\)-module \(M''\) is torsion-free of finite type, and consequently free of rank \textless{} d.~Arguing by induction on d and applying the induction hypothesis to the endomorphism of \(M''\) induced by u, we conclude that all the eigenvalues of v in an algebraically closed extension of \(\mathbf{Q}\) are integers.

\begin{enumerate}
\def\labelenumi{\alph{enumi})}
\setcounter{enumi}{2}
\tightlist
\item
  We show that \(v\) is diagonalizable. Let \(\lambda\) be an eigenvalue of \(v\) and let \(x \in \mathrm{M} \otimes_{\mathbf{Z}} \mathbf{Q}\) be such that \((v - \lambda)^2 x = 0\) . We have \(v(vx - \lambda x) = \lambda(vx - \lambda x)\) , so
\end{enumerate}

\[
\begin{array}{c} \frac {1}{n !} (v - \lambda - n + 1) (v - \lambda - n + 2) \ldots (v - \lambda - 1) (v - \lambda) x \\ = \frac {(- 1) ^ {n - 1}}{n} (v x - \lambda x). \end{array}
\]

By \(a\) ), this implies that \(vx - \lambda x \in n\mathbf{M}\) for all \(n \in \mathbf{N}\) , so \((v - \lambda)x = 0\) .

\begin{enumerate}
\def\labelenumi{\alph{enumi})}
\setcounter{enumi}{3}
\tightlist
\item
  Let \(\lambda\) be an eigenvalue of v and let \((\lambda - a, \lambda + b]\) be an interval in \(\mathbf{Z}\) containing all the eigenvalues of v. Consider the polynomial
\end{enumerate}

\[
\begin{array}{c} \mathrm{P(T)} = (- 1) ^ {b} \frac {(\mathrm{T} - \lambda - 1) (\mathrm{T} - \lambda - 2) \ldots (\mathrm{T} - \lambda - b)}{b !} \\ \times \frac {(\mathrm{T} - \lambda + 1) (\mathrm{T} - \lambda + 2) \ldots (\mathrm{T} - \lambda + a)}{a !}. \end{array}
\]

We have \(\mathrm{P}(\mathbf{Z})\subset\mathbf{Z},\mathrm{P}(\lambda)=1,\mathrm{P}(\mu)=0\) for \(\mu\in\mathbf{Z}\cap(\lambda-a,\lambda+b)\) and \(\mu\neq\lambda\) . By a), \(\mathrm{P}(v)(\mathrm{M})\subset\mathrm{M}\) . By c), \(\mathrm{P}(v)\) is a projection of \(M\otimes_{\mathbf{Z}}\mathbf{Q}\) onto the eigenspace corresponding to \(\lambda\) .

Q.E.D.

Remark 1. If we only assume that v is diagonalizable with integer eigenvalues, u is not necessarily diagonalizable (for example, take \(M = \mathbf{Z}^{2}\) and \(u(x, y) = (y, x)\) for all \((x, y) \in M\) ).

Let \(\mathfrak{g},\mathfrak{h},\mathrm{R},\mathcal{H},\mathcal{G}^{\alpha},\mathcal{G},\mathcal{U}\) be as in no. 7, and assume that \(\mathcal{G}\) is a Chevalley order in \((\mathfrak{g},\mathfrak{h})\) .

DEFINITION 3. Let E be a \(\mathfrak{g}\)-module. A lattice \(\mathcal{E}\) in E is said to be admissible (relative to \(\mathcal{G}\)) if the following conditions are satisfied:

\begin{enumerate}
\def\labelenumi{(\roman{enumi})}
\item
  \(\mathcal{U}\) maps \(\mathcal{E}\) to \(\mathcal{E}\) ;
\item
  \(\mathcal{E}\) is stable under \(\binom{h}{n}\) and \(x^{(n)}\) for all \(\alpha \in \mathrm{R}, x \in \mathcal{G}^{\alpha}, n \in \mathbf{N}, h \in \mathcal{H}\) .
\end{enumerate}

Remarks. 2) Let \(\rho\) be the adjoint representation of \(\mathfrak{g}\) on \(\mathrm{U}(\mathfrak{g})\) . Let \(\alpha, x, n, h\) be as in (ii) above. We have \(\rho(x^{(n)})\) . \(\mathcal{U} \subset \mathcal{U}\) by Lemma 2. On the other hand, if \(p \in \mathbf{N}\) ,

\[
\rho \left(\binom{h}{p}\right) x ^ {(n)} = \binom{\operatorname{ad} h}{p} x ^ {(n)} = \binom{n \alpha (h)}{p} x ^ {(n)}
\]

(no. 5, formula (13)), so \(\rho\left(\binom{h}{p}\right). \mathcal{U} \subset \mathcal{U}\) . This proves that \(\mathcal{U}\) is an admissible lattice in \(\mathrm{U}(\mathfrak{g})\) , and it follows that \(\mathcal{G}\) is an admissible lattice in \(\mathfrak{g}\) (for the adjoint representation).

\begin{enumerate}
\def\labelenumi{\arabic{enumi})}
\setcounter{enumi}{2}
\tightlist
\item
  Let E be a finite dimensional \(\mathfrak{g}\)-module, \(\mathcal{E}\) an admissible lattice in E, \(\mathfrak{c}\) the centre of \(\mathfrak{g}\). By Lemma 6, every element of \(\mathfrak{c}\) defines a diagonalizable endomorphism of E. Hence E is semi-simple (Chap. I, §6, no. 5, Th. 4). Thus, E is a direct sum of simple \(\mathcal{D}\mathfrak{g}\)-modules on which \(\mathfrak{c}\) induces homotheties. By Lemma 6, \(\mathcal{E}=\oplus(\mathcal{E}\cap\mathrm{E}^{\lambda})\) and, for all weights \(\lambda\) of E, we have
\end{enumerate}

\[
\lambda (\mathcal {H}) \subset \mathbf {Z}.
\]

\begin{enumerate}
\def\labelenumi{\arabic{enumi})}
\setcounter{enumi}{3}
\tightlist
\item
  If \(\mathfrak{g}\) is semi-simple and \(\mathcal{H} = \mathrm{Q}(\mathrm{R}^{\vee})\) , conditions (i) and (ii) of Def. 3 are equivalent, by Th. 3 (ii), to
\end{enumerate}

\begin{enumerate}
\def\labelenumi{(\roman{enumi})}
\setcounter{enumi}{2}
\tightlist
\item
  \(\mathcal{E}\) is stable under \(x^{(n)}\) for all \(\alpha \in \mathrm{R}, x \in \mathcal{G}^{\alpha}, n \in \mathbf{N}\) .
\end{enumerate}

\begin{enumerate}
\def\labelenumi{\arabic{enumi})}
\setcounter{enumi}{4}
\tightlist
\item
  Let B be a basis of R; in conditions (i) and (ii) above, `` \(\alpha \in R\) '' can be replaced by `` \(\alpha \in B \cup (-B)\) '' (loc. cit).
\end{enumerate}

THEOREM 4. Let \(\mathbf{E}\) be a finite dimensional \(\mathfrak{g}\) -module. The following conditions are equivalent:

\begin{enumerate}
\def\labelenumi{(\roman{enumi})}
\item
  E has an admissible lattice;
\item
  every element of \(\mathcal{H}\) defines a diagonalizable endomorphism of \(E\) with integer eigenvalues.
\item
  \(\Longrightarrow\) (ii): this follows from Remark 3.
\item
  \(\Longrightarrow\) (i): we assume that condition (ii) is satisfied and prove (i). By Th. 4 of Chap. I, §6, no. 5, we can assume that the elements of \(\mathfrak{c}\) define homotheties of E, and that E is a simple \(\mathcal{D}\mathfrak{g}\)-module. Let B be a basis of \(\mathrm{R}\), and \(\mathfrak{g} = \mathfrak{n}_{-} \oplus \mathfrak{h} \oplus \mathfrak{n}_{+}\) the corresponding decomposition of \(\mathfrak{g}\). Let \(\lambda\) be the highest weight of the \(\mathcal{D}\mathfrak{g}\)-module E, and let \(e \in E^{\lambda} - \{0\}\) . Put \(\mathcal{E} = \mathcal{U}.e\). It is clear that \(\mathcal{U}.\mathcal{E} \subset \mathcal{E}\) . Since E is simple, \(\mathrm{U}(\mathfrak{g})\) . e = E and hence \(\mathcal{E}\) generates E as a \(\mathbf{Q}\)-vector space. For \(h \in \mathcal{H}\) and \(n \in \mathbf{N}\) , we have \(\binom{h}{n} e = \binom{\lambda(h)}{n} e \in \mathbf{Z} e\) , so
\end{enumerate}

\[
(\mathcal {U} \cap \mathrm{U} (\mathfrak {h})). e = \mathbf {Z}   e.
\]

Since \(\mathrm{U}(\mathfrak{n}_{+}).e=0\) , we have \(\mathcal{E}=(\mathcal{U}\cap\mathrm{U}(\mathfrak{n}_{-}))\) .e by Prop. 3. It now follows from Th. 3 (iii) that \(\mathcal{E}\) is a \(\mathbf{Z}\)-module of finite type.

COROLLARY. If \(\mathfrak{g}\) is semi-simple and \(\mathcal{H}=\mathrm{Q}(\mathrm{R}^{\vee})\) , every finite dimensional \(\mathfrak{g}\)-module has an admissible lattice.

\section{CLASSICAL SPLITTABLE SIMPLE LIE ALGEBRAS}

In this paragraph we describe explicitly, for each type of classical splittable simple Lie algebra:

\begin{enumerate}
\def\labelenumi{(\Roman{enumi})}
\item
  an algebra of this type, its dimension and its splitting Cartan subalgebras;
\item
  its coroots;
\item
  its Borel subalgebras and its parabolic subalgebras;
\item
  its fundamental simple representations;
\item
  those of its fundamental simple representations which are orthogonal or symplectic;
\item
  the algebra of invariant polynomial functions;
\item
  certain properties of the groups \(\operatorname{Aut} \mathfrak{g}\) , \(\operatorname{Aut}_0 \mathfrak{g}\) , \(\operatorname{Aut}_e \mathfrak{g}\) ;
\item
  the restriction of the Killing form to a Cartan subalgebra;
\item
  the Chevalley orders.
\end{enumerate}

\subsection{\texorpdfstring{ALGEBRAS OF TYPE \(\mathbf{A}_l\) ( \(l \geq 1\) )}{ALGEBRAS OF TYPE \textbackslash mathbf\{A\}\_l ( l \textbackslash geq 1 )}}

\begin{enumerate}
\def\labelenumi{(\Roman{enumi})}
\tightlist
\item
  Let V be a vector space of dimension \(l+1\) over k, and let \(\mathfrak{g}\) be the algebra \(\mathfrak{sl}(V)\) of endomorphisms of V of trace zero. Let \((e_{i})_{1\leq i\leq l+1}\) be a basis of V; the map which associates to an element of \(\mathfrak{g}\) its matrix with respect to this basis is an identification of \(\mathfrak{g}\) with the algebra \(\mathfrak{sl}(l+1,k)\) of matrices of trace zero. We know that \(\mathfrak{g}\) is semi-simple (Chap. I, §6, no. 7, Prop. 8).
\end{enumerate}

Recall (Algebra, Chap. II, §10, no. 3) that \(E_{ij}\) denotes the matrix \((\alpha_{mp})\) such that \(\alpha_{ij} = 1\) and \(\alpha_{mp} = 0\) for \((m,p) \neq (i,j)\) . The matrices

\[
\begin{array}{l l} E _ {i j} & (1 \leq i, j \leq l + 1, i \neq j) \\ E _ {i, i} - E _ {i + 1, i + 1} & (1 \leq i \leq l) \end{array}
\]

form a basis of \(\mathfrak{g}\). Hence

\[
\dim \mathfrak {g} = l (l + 2).
\]

Let \(\hat{\mathfrak{h}}\) be the set of diagonal elements of \(\mathfrak{gl}(l + 1,k)\) ; the sequence \((E_{ii})_{1\leq i\leq l + 1}\) is a basis of the vector space \(\hat{\mathfrak{h}}\) ; let \((\hat{\varepsilon}_i)_{1\leq i\leq l + 1}\) be the basis of \(\hat{\mathfrak{h}}^*\) dual to \((E_{ii})_{1\leq i\leq l + 1}\) . For all \(h\in \hat{\mathfrak{h}}\) ,

\[
[ h, E _ {i j} ] = (\hat {\varepsilon} _ {i} (h) - \hat {\varepsilon} _ {j} (h)) E _ {i j}\tag{1}
\]

by Chap. I, §1, no. 2, formulas (5). Let \(\mathfrak{h}\) be the set of elements of \(\hat{\mathfrak{h}}\) of trace zero, and put \(\varepsilon_i = \hat{\varepsilon}_i|\mathfrak{h}\) . Then \(\mathfrak{h}\) is a Cartan subalgebra of \(\mathfrak{g}\) (Chap. VII, §2, no. 1, Example 4). Relation (1) proves that this Cartan subalgebra is splitting, and that the roots of \((\mathfrak{g},\mathfrak{h})\) are the \(\varepsilon_i - \varepsilon_j\) ( \(i\neq j\) ). Let \(\hat{\mathfrak{h}}_0^*\) be the set of elements of \(\hat{\mathfrak{h}}^*\) the sum of whose coordinates with respect to \((\hat{\varepsilon}_i)\) is zero. The map \(\lambda \mapsto \lambda |\mathfrak{h}\) from \(\hat{\mathfrak{h}}_0^*\) to \(\mathfrak{h}^*\) is bijective. Thus, the root system R of \((\mathfrak{g},\mathfrak{h})\) is of type \(A_l\) (Chap. VI, §4, no. 7). Consequently, \(\mathfrak{g}\) is simple (§3, no. 2, Cor. 1 of Prop. 6). Thus, \(\mathfrak{g}\) is a splittable simple Lie algebra of type \(A_l\) .

Every splitting Cartan subalgebra \(\mathfrak{h}^{\prime}\) of \(\mathfrak{g}\) is a transform of \(\mathfrak{h}\) under an elementary automorphism ( \(\S3\) , no. 3, Cor. of Prop. 10). Since \(\operatorname{Aut}_{e}\mathfrak{g}\) is the set of automorphisms \(x \mapsto sxs^{-1}\) of \(\mathfrak{g}\) with \(s \in \mathbf{SL}(V)\) (Chap. VII, \(\S3\) , no. 1, Remark 2; cf.~also (VII)), there exists a basis \(\beta\) of V such that \(\mathfrak{h}^{\prime}\) is the set \(\mathfrak{h}_{\beta}\) of elements of \(\mathfrak{g}\) whose matrix with respect to the basis \(\beta\) is diagonal. Since \(\mathfrak{h}_{\beta}\) contains an element with distinct eigenvalues, the only vector subspaces of V stable under the elements of \(\mathfrak{h}_{\beta}\) are those generated by a subset of \(\beta\) . It follows that the map \(\beta \mapsto \mathfrak{h}_{\beta}\) induces by passage to the quotient a bijection from the set of decompositions of V into the direct sum of \(l + 1\) subspaces of dimension 1 to the set of splitting Cartan subalgebras of \(\mathfrak{g}\).

\begin{enumerate}
\def\labelenumi{(\Roman{enumi})}
\setcounter{enumi}{1}
\tightlist
\item
  Let \(\alpha = \varepsilon_{i} - \varepsilon_{j}\) ( \(i \neq j\) ) be a root. We have \(\mathfrak{g}^{\alpha} = kE_{ij}\) . Since
\end{enumerate}

\[
[ E _ {i j}, E _ {j i} ] = E _ {i i} - E _ {j j}
\]

and since \(\alpha(E_{ii} - E_{jj}) = 2\) , we have (§2, no. 2, Th. 1 (ii))

\[
H _ {\alpha} = E _ {i i} - E _ {j j}.
\]

\begin{enumerate}
\def\labelenumi{(\Roman{enumi})}
\setcounter{enumi}{2}
\tightlist
\item
  Put \(\alpha_{1} = \varepsilon_{1} - \varepsilon_{2},\alpha_{2} = \varepsilon_{2} - \varepsilon_{3},\ldots ,\alpha_{l} = \varepsilon_{l} - \varepsilon_{l + 1}\) . By Chap. VI, §4, no. 7.I, \((\alpha_{1},\dots,\alpha_{l})\) is a basis B of R; the positive roots relative to B are the \(\varepsilon_i - \varepsilon_j\) for \(i < j\) . The corresponding Borel subalgebra \(\mathfrak{b}\) is the set of upper triangular matrices of trace zero.
\end{enumerate}

A flag in V is a set of vector subspaces of V, distinct from \(\{0\}\) and V, totally ordered by inclusion. Order the set of flags of V by inclusion. The maximal flags are the sets \(\{W_{1},\ldots,W_{l}\}\) , where \(W_{i}\) is an i-dimensional vector subspace and

\[
\mathrm{W} _ {1} \subset \dots \subset \mathrm{W} _ {l}.
\]

For example, if \(V_{i}\) denotes the subspace of V generated by \(e_{1},\ldots,e_{i}\) , then \(\{V_{1},\ldots,V_{l}\}\) is a maximal flag.

It is immediate that \(\mathfrak{b}\) is the set of elements of \(\mathfrak{g}\) leaving stable the elements of the maximal flag \(\{V_{1},\ldots,V_{l}\}\) . Conversely, since \(\mathfrak{b}\) contains \(\mathfrak{h}\) and the matrices \(E_{ij}\) for i\textless j, we see that the \(V_{i}\) are the only non-trivial vector subspaces stable under \(\mathfrak{b}\).

Now let \(\delta\) be a maximal flag in V. It follows from the preceding that the set \(\mathfrak{b}_{\delta}\) of elements of \(\mathfrak{g}\) leaving stable all the elements of \(\delta\) is a Borel subalgebra of \(\mathfrak{g}\). Since every Borel subalgebra of \(\mathfrak{g}\) is a transform of \(\mathfrak{b}\) under an elementary automorphism, we see that the map \(\delta \mapsto \mathfrak{b}_{\delta}\) is a bijection from the set of maximal flags to the set of Borel subalgebras of \(\mathfrak{g}\).

Let \(\beta\) be a basis of V. By (I) and the preceding, the Borel subalgebras containing \(\mathfrak{h}_{\beta}\) are those corresponding to the maximal flags each of whose elements is generated by a subset of \(\beta\) . These flags correspond bijectively to the total orders on \(\beta\) in the following way: to a total order \(\omega\) on \(\beta\) is associated the flag \(\{\mathrm{W}_1,\dots ,\mathrm{W}_l\}\) , where \(\mathrm{W}_i\) is the vector subspace generated by the first \(i\) elements of \(\beta\) for the order \(\omega\) . Since there are \((l + 1)!\) total orders on \(\beta\) , we recover the fact that there exist \((l + 1)!\) Borel subalgebras of \((\mathfrak{sl}(\mathrm{V}),\mathfrak{h}_{\beta})\) ( \(\S 3\) , no. 3, Remark).

Let \(\gamma\) be a flag in V. Since \(\gamma\) is contained in a maximal flag, the set \(\mathfrak{p}_{\gamma}\) of elements of \(\mathfrak{g}\) leaving stable the elements of \(\gamma\) is a parabolic subalgebra of \(\mathfrak{g}\). We show that the only non-trivial vector subspaces stable under \(\mathfrak{p}_{\gamma}\) are the elements of \(\gamma\) . For this, we can assume that \(\gamma = \{V_{i_1}, \ldots, V_{i_q}\}\) with \(1 \leq i_1 < \cdots < i_q \leq l\) . Put \(i_0 = 0, i_{q+1} = l + 1\) . The non-empty intervals

\[
\left[ i _ {0} + 1, i _ {1} \right], \left[ i _ {1} + 1, i _ {2} \right], \dots , \left[ i _ {q} + 1, i _ {q + 1} \right]
\]

form a partition of \(\{1, \ldots, l + 1\}\) , so that any square matrix of order \(l + 1\) can be written as a block matrix \((X_{ab})_{1 \leq a, b \leq q + 1}\) . The algebra \(\mathfrak{p}_{\gamma}\) is then the set \(\mathfrak{p}_{i_1,\dots,i_q}\) of elements \((X_{ab})_{1\leq a,b\leq q + 1}\) of \(\mathfrak{sl}(l + 1,k)\) such that \(X_{ab} = 0\) for \(a > b\) . Since \(\mathfrak{p}_{i_1,\dots,i_q}\supset \mathfrak{b}\) , a non-trivial vector subspace stable under \(\mathfrak{p}_{i_1,\dots,i_q}\) is one of the \(\mathrm{V}_i\) ; if \(i_k < i < i_{k + 1}\) , the algebra \(\mathfrak{p}_{i_1,\dots,i_q}\) contains \(E_{i_{k + 1},i}\) and \(\mathrm{V}_i\) is not stable, hence our assertion.

Consequently, the \(2^{l}\) flags contained in the maximal flag \(\{V_{1},\ldots,V_{l}\}\) give rise to \(2^{l}\) distinct parabolic subalgebras containing \(\mathfrak{b}\); since there are exactly \(2^{l}\) parabolic subalgebras containing \(\mathfrak{b}\) ( \(\S3\) , no. 4, Remark), it follows that the map \(\gamma\mapsto \mathfrak{p}_{\gamma}\) is a bijection from the set of flags of V to the set of parabolic subalgebras of \(\mathfrak{g}\). Moreover, \(\mathfrak{p}_{\gamma}\supset \mathfrak{p}_{\gamma'}\) if and only if \(\gamma\subset\gamma'\) .

Recall the parabolic subalgebra \(\mathfrak{p} = \mathfrak{p}_{i_1, \ldots, i_q}\) ( \(1 \leq i_1 < \cdots < i_q \leq l\) ). Let \(\mathfrak{s}\) (resp. \(\mathfrak{n}\)) be the set of \((X_{ab})_{1 \leq a, b \leq q+1}\) in \(\mathfrak{sl}(l+1, k)\) such that \(X_{ab} = 0\) for \(a \neq b\) (resp. \(a \geq b\) ). In view of Prop. 13 of §3, no. 4, we have \(\mathfrak{p} = \mathfrak{s} \oplus \mathfrak{n}\) , the subalgebra \(\mathfrak{s}\) is reductive in \(\mathfrak{g}\) and \(\mathfrak{n}\) is both the largest nilpotent ideal and the nilpotent radical of \(\mathfrak{p}\).

\begin{enumerate}
\def\labelenumi{(\Roman{enumi})}
\setcounter{enumi}{3}
\tightlist
\item
  For \(r = 1, 2, \ldots, l\) , let \(\varpi_{r} = \varepsilon_{1} + \cdots + \varepsilon_{r}\) . We have \(\varpi_{i}(H_{\alpha_{j}}) = \delta_{ij}\) , so \(\varpi_{r}\) is the fundamental weight corresponding to \(\alpha_{r}\) .
\end{enumerate}

Let \(\sigma\) be the identity representation of \(\mathfrak{g}\) on V. The exterior power \(\bigwedge^{r}\sigma\) of \(\sigma\) is a representation on \(\mathrm{E} = \bigwedge^{r}(\mathrm{V})\) . Let \((e_1,\ldots ,e_{l + 1})\) be the chosen basis of V. The \(e_{i_1}\wedge \dots \wedge e_{i_r}\) , where \(i_1 < \dots < i_r\) , form a basis of E. If \(h\in \mathfrak{h}\) ,

\[
(\bigwedge^ {r} \sigma) (h). e _ {i _ {1}} \wedge \dots \wedge e _ {i _ {r}} = (\varepsilon_ {i _ {1}} + \dots + \varepsilon_ {i _ {r}}) (h) e _ {i _ {1}} \wedge \dots \wedge e _ {i _ {r}}.
\]

Thus, every weight is of multiplicity 1, \(\varpi_{r}\) is a weight of \(\bigwedge^{r}\sigma\) , and every other weight is of the form \(\varpi_{r}-\mu\) , where \(\mu\) is a positive radical weight. Consequently, \(\varpi_{r}\) is the highest weight of \(\bigwedge^{r}\sigma\) , and \(e_{1}\wedge\cdots\wedge e_{r}\) is a primitive element. By Chap. VI, §4, no. 7.IX, the Weyl group can be identified with the symmetric group of

\[
\{\varepsilon_ {1}, \dots , \varepsilon_ {l + 1} \}.
\]

The orbit of \(\varpi_{r}\) under the Weyl group thus contains all the \(\varepsilon_{i_{1}} + \cdots + \varepsilon_{i_{r}}\) with \(i_{1} < \cdots < i_{r}\) . The simple submodule generated by the primitive element \(e_{1} \wedge \cdots \wedge e_{r}\) thus admits all the \(\varepsilon_{i_{1}} + \cdots + \varepsilon_{i_{r}}\) as weights and consequently is equal to E. Thus, \(\bigwedge^{r} \sigma\) is irreducible with highest weight \(\varpi_{r}\) .

Thus, the representations \(\bigwedge^{r}\sigma\) ( \(1 \leq r \leq l\) ) are the fundamental representations. We have \(\dim(\bigwedge^{r}\sigma) = \binom{l+1}{r}\) .

\begin{enumerate}
\def\labelenumi{(\Alph{enumi})}
\setcounter{enumi}{21}
\tightlist
\item
  We have \(w_0(\alpha_1) = -\alpha_l, w_0(\alpha_2) = -\alpha_{l-1}, \ldots\) (Chap. VI, §4, no. 7, XI), so
\end{enumerate}

\[
- w _ {0} (\varpi_ {1}) = \varpi_ {l}, - w _ {0} (\varpi_ {2}) = \varpi_ {l - 1}, \ldots .
\]

Let

\[
\omega = n _ {1} \varpi_ {1} + \dots + n _ {l} \varpi_ {l} \quad (n _ {1}, \ldots , n _ {l} \in \mathbf {N})
\]

be a dominant weight. Then, the simple representation with highest weight \(\omega\) is orthogonal or symplectic if and only if

\[
n _ {1} = n _ {l}, \quad n _ {2} = n _ {l - 1}, \ldots
\]

(§7, no. 5, Prop. 12). In particular, if \(l\) is even, none of the fundamental representations of \(\mathfrak{sl}(l + 1,k)\) is orthogonal or symplectic. If \(l\) is odd, the representation \(\bigwedge^i\sigma\) for \(i\neq (l + 1) / 2\) is neither orthogonal nor symplectic; by Chap. VI, §4, no. 7.VI, the sum of the coordinates of \(\varpi_{(l + 1) / 2}\) with respect to \((\alpha_{1},\ldots ,\alpha_{l})\) is

\[
\begin{array}{r l} \frac {1}{l + 1} \left[ \frac {l + 1}{2} \left(1 + 2 + \dots + \frac {l - 1}{2}\right) + \frac {l + 1}{2} \left(1 + 2 + \dots + \frac {l + 1}{2}\right) \right] & = 1 + 2 + \dots + \frac {l - 1}{2} + \frac {l + 1}{4} \end{array}
\]

so \(\bigwedge^{(l + 1) / 2}\sigma\) is orthogonal if \(l\equiv -1\) (mod. 4) and symplectic if \(l\equiv 1\) (mod. 4) (§7, no. 5, Prop. 12). This last result can be made more precise as follows. Choose a non-zero element \(e\) in \(\bigwedge^{l + 1}(\mathrm{V})\) . The multiplication in the exterior algebra \(\mathrm{V}\) defines a bilinear map from

\[
\bigwedge^ {(l + 1) / 2} (\mathrm{V}) \times \bigwedge^ {(l + 1) / 2} (\mathrm{V})
\]

to \(\bigwedge^{l + 1}(\mathrm{V})\) , which can be written \((u,v)\mapsto \varPhi (u,v)e\) , where \(\varPhi\) is a bilinear form on \(\bigwedge^{(l + 1) / 2}(\mathrm{V})\) . It is immediately verified that \(\varPhi\) is non-zero, invariant under \(\mathfrak{g}\) (and hence non-degenerate), symmetric if \((l + 1) / 2\) is even, and alternating if \((l + 1) / 2\) is odd.

\begin{enumerate}
\def\labelenumi{(\Roman{enumi})}
\setcounter{enumi}{5}
\tightlist
\item
  For all \(x \in \mathfrak{g}\) , the characteristic polynomial of \(\sigma(x) = x\) can be written
\end{enumerate}

\[
\mathrm{T} ^ {l + 1} + f _ {2} (x) \mathrm{T} ^ {l - 1} + f _ {3} (x) \mathrm{T} ^ {l - 2} + \dots + f _ {l + 1} (x)
\]

where \(f_{2}, \ldots, f_{l+1}\) are polynomial functions invariant under \(\mathfrak{g}\) (§8, no. 3, Lemma 2).

If \(x = \xi_{1}E_{11} + \cdots + \xi_{l+1}E_{l+1 l+1} \in \mathfrak{h}\) , the \(f_{i}(x)\) are, up to sign, the elementary symmetric functions of \(\xi_{1}, \ldots, \xi_{l+1}\) of degree \(2, \ldots, l + 1\) . Thus, by Chap. VI, §4, no. 7.IX, the \(f_{i}|_{\mathfrak{h}}\) generate the algebra of elements of \(\mathbf{S}(\mathfrak{h}^{*})\) invariant under the Weyl group, and are algebraically independent. Hence (§8, no. 3, Prop. 3) \(f_{2}, f_{3}, \ldots, f_{l+1}\) generate the algebra of polynomial functions invariant under \(\mathfrak{g}\), and are algebraically independent.

\begin{enumerate}
\def\labelenumi{(\Roman{enumi})}
\setcounter{enumi}{6}
\tightlist
\item
  For all \(g \in \mathbf{GL}(l + 1, k)\) , let \(\varphi_k(g) = \varphi(g)\) be the automorphism \(x \mapsto gxg^{-1}\) of \(\mathfrak{g}\) . Then \(\varphi\) is a homomorphism from \(\mathbf{GL}(l + 1, k)\) to \(\operatorname{Aut}(\mathfrak{g})\) . We have
\end{enumerate}

\[
\varphi (\mathbf {S L} (l + 1, k)) = \operatorname{Aut} _ {e} (\mathfrak {g})
\]

(Chap. VII, §3, no. 1, Remark 2). Let \(\bar{k}\) be an algebraic closure of \(k\) . We have

\[
\mathbf {G L} (l + 1, \bar {k}) = \bar {k} ^ {*}. \mathbf {S L} (l + 1, \bar {k}),
\]

so \(\varphi_{\bar{k}}(\mathbf{GL}(l + 1,\bar{k})) = \varphi_{\bar{k}}(\mathbf{SL}(l + 1,\bar{k})) = \mathrm{Aut}_e(\mathfrak{g}\otimes_k\bar{k})\) ; it follows that \(\varphi (\mathbf{GL}(l + 1,k))\subset \mathrm{Aut}_0(\mathfrak{g})\) . On the other hand, \(\mathrm{Aut}_0(\mathfrak{g})\subset \varphi (\mathbf{GL}(l + 1,k))\) , by Prop. 2 of §7, no. 1, applied to the identity representation of \(\mathfrak{g}\) . Hence,

\[
\operatorname{Aut} _ {0} (\mathfrak {g}) = \varphi (\mathbf {G L} (l + 1, k)).
\]

The kernel of \(\varphi\) is the set of elements of \(\mathbf{GL}(l+1,k)\) that commute with every matrix of order \(l+1\) , that is the set \(k^{*}\) of invertible scalar matrices. Thus, \(\operatorname{Aut}_{0}(\mathfrak{g})\) can be identified with the group \(\mathbf{GL}(l+1,k)/k^{*}=\mathbf{PGL}(l+1,k)\) . The kernel of \(\varphi'=\varphi|\mathbf{SL}(l+1,k)\) is \(\mu_{l+1}(k)\) , where \(\mu_{l+1}(k)\) denotes the set of \((l+1)\) th roots of unity in k. Thus, \(\operatorname{Aut}_{e}(\mathfrak{g})\) can be identified with the group \(\mathbf{SL}(l+1,k)/\mu_{l+1}(k)=\mathbf{PSL}(l+1,k)\) . On the other hand, we have the exact sequence

\[
1 \longrightarrow \mathbf {S L} (l + 1, k) \longrightarrow \mathbf {G L} (l + 1, k) \stackrel {\det} {\longrightarrow} k ^ {*} \longrightarrow 1
\]

and the image of \(k^*\) under det is \(k^{*l + 1}\) . It follows that there are canonical isomorphisms

\[
\begin{array}{c} \operatorname{Aut} _ {0} (\mathfrak {g}) / \operatorname{Aut} _ {e} (\mathfrak {g}) \longrightarrow \mathbf {P G L} (l + 1, k) / \mathbf {P S L} (l + 1, k) \\ \qquad \qquad \qquad \qquad \qquad \longrightarrow \mathbf {G L} (l + 1, k) / k ^ {*}. \mathbf {S L} (l + 1, k) \longrightarrow k ^ {*} / k ^ {*   l + 1}. \end{array}
\]

If \(k = \mathbf{R}\) , we see that \(\mathrm{Aut}_0(\mathfrak{g}) = \mathrm{Aut}_e(\mathfrak{g})\) if \(l + 1\) is odd, and that \(\mathrm{Aut}_0(\mathfrak{g}) / \mathrm{Aut}_e(\mathfrak{g})\) is isomorphic to \(\mathbf{Z}/2\mathbf{Z}\) if \(l + 1\) is even.

With the notations of §5, \(f(\mathrm{T}_{\mathbf{Q}})\) is the set of automorphisms of \(\mathfrak{g}\) that induce the identity on \(\mathfrak{h}\), and hence is equal to \(\varphi(\mathrm{D})\) , where D is the set of diagonal elements of \(\mathbf{GL}(l+1,k)\) (§5, Prop. 4). Let \(D'\) be the set of diagonal elements of \(\mathbf{SL}(l+1,k)\) . By Prop. 3 of §5, and the determination of \(\operatorname{Aut}_{e}(\mathfrak{g})\) , we have \(f(q(\mathrm{T}_{\mathrm{P}})) \subset \varphi(\mathrm{D}')\) . We show that \(f(q(\mathrm{T}_{\mathrm{P}})) = \varphi(\mathrm{D}')\) . Let

\[
d = \left( \begin{array}{c c c} \lambda_ {1} & & 0 \\ & \ddots & \\ 0 & & \lambda_ {l + 1} \end{array} \right)
\]

be an element of \(D'\) . There exists a \(\zeta \in \operatorname{Hom}(\mathbf{Q}(\mathbf{R}), k^{*}) = \mathrm{T}_{\mathbf{Q}}\) such that \(\zeta(\varepsilon_{i} - \varepsilon_{j}) = \lambda_{i}\lambda_{j}^{-1}\) for all i and j. It is easy to verify that \(f(\zeta) = \varphi(d)\) . By Chap. VI, §4, no. 7.VIII, \(\mathrm{P}(\mathrm{R})\) is generated by \(\mathrm{Q}(\mathrm{R})\) and the element \(\varepsilon = \varepsilon_{1}\) , whose image in \(\mathrm{P}(\mathrm{R})/\mathrm{Q}(\mathrm{R})\) is of order \(l + 1\) ; but

\[
\begin{array}{c} \zeta ((l + 1) \varepsilon) = \zeta ((\varepsilon_ {1} - \varepsilon_ {2}) + (\varepsilon_ {1} - \varepsilon_ {3}) + \dots + (\varepsilon_ {1} - \varepsilon_ {l + 1})) \\ = \lambda_ {1} ^ {l} \lambda_ {2} ^ {- 1} \lambda_ {3} ^ {- 1} \ldots \lambda_ {l + 1} ^ {- 1} = \lambda_ {1} ^ {l + 1} \end{array}
\]

so \(\zeta\) extends to a homomorphism from \(\mathrm{P}(\mathrm{R})\) to \(k^{*}\) . This proves that \(\zeta \in q(\mathrm{T}_{\mathrm{P}})\) , so \(\varphi(d) \in f(q(\mathrm{T}_{\mathrm{P}}))\) .

Recall (§5, no. 3, Cor. 2 of Prop. 5) that \(\mathrm{Aut}(\mathfrak{g}) = \mathrm{Aut}_0(\mathfrak{g})\) for \(l = 1\) , and that \(\mathrm{Aut}(\mathfrak{g}) / \mathrm{Aut}_0(\mathfrak{g})\) is isomorphic to \(\mathbf{Z}/2\mathbf{Z}\) for \(l \geq 2\) . The map \(\theta : x \mapsto -^t x\) is an automorphism of \(\mathfrak{sl}(l + 1, k)\) and \(a_0 = \theta|\mathfrak{h} \notin \mathrm{W}\) if \(l \geq 2\) (Chap. VI, §4, no. 7.XI), so the class of \(a_0\) in \(\mathrm{Aut}(\mathfrak{g}) / \mathrm{Aut}_0(\mathfrak{g})\) is the non-trivial element of this group (§5, no. 2, Prop. 4).

\begin{enumerate}
\def\labelenumi{(\Roman{enumi})}
\setcounter{enumi}{7}
\tightlist
\item
  The restriction to \(\mathfrak{h}\) of the Killing form is
\end{enumerate}

\[
\begin{array}{l} \Phi (\xi_ {1} E _ {1 1} + \dots + \xi_ {l + 1} E _ {l + 1, l + 1}, \xi_ {1} ^ {\prime} E _ {1 1} + \dots + \xi_ {l + 1} ^ {\prime} E _ {l + 1, l + 1}) \\ = \sum_ {i \neq j} (\xi_ {i} - \xi_ {j}) (\xi_ {i} ^ {\prime} - \xi_ {j} ^ {\prime}) = \sum_ {i, j} (\xi_ {i} - \xi_ {j}) (\xi_ {i} ^ {\prime} - \xi_ {j} ^ {\prime}) \\ = (l + 1) \sum_ {i} \xi_ {i} \xi_ {i} ^ {\prime} + (l + 1) \sum_ {j} \xi_ {j} \xi_ {j} ^ {\prime} - 2 \left(\sum_ {i} \xi_ {i}\right) \left(\sum_ {j} \xi_ {j} ^ {\prime}\right) \\ = 2 (l + 1) \sum_ {i} \xi_ {i} \xi_ {i} ^ {\prime}. \end{array}
\]

\begin{enumerate}
\def\labelenumi{(\Roman{enumi})}
\setcounter{enumi}{8}
\tightlist
\item
  For \(1 \leq i < j \leq l + 1\) , put
\end{enumerate}

\[
X _ {\varepsilon_ {i} - \varepsilon_ {j}} = E _ {i j} X _ {\varepsilon_ {j} - \varepsilon_ {i}} = - E _ {j i}.
\]

Then, for all \(\alpha \in \mathrm{R}\) , we have \([X_{\alpha}, X_{-\alpha}] = -H_{\alpha}\) and \(\theta(X_{\alpha}) = X_{-\alpha}\) (where \(\theta\) is the automorphism \(x \mapsto -^{t} x\) introduced in (VII)). Consequently, \((X_{\alpha})_{\alpha \in \mathrm{R}}\) is a Chevalley system in \((\mathfrak{g}, \mathfrak{h})\) .

Take \(k = \mathbf{Q}\) . The permissible lattices in \(\mathfrak{h}\) (§12, no. 6, Def. 1) are those lying between the \(\mathbf{Z}\) -module \(\mathrm{Q}(\mathrm{R}^{\vee})\) generated by the \(E_{ii} - E_{i+1,i+1}\) , that is, consisting of the diagonal matrices belonging to \(\mathfrak{sl}(l+1,\mathbf{Z})\) , and the \(\mathbf{Z}\) -module \(\mathrm{P}(\mathrm{R}^{\vee})\) generated by \(\mathrm{Q}(\mathrm{R}^{\vee})\) and \(E_{11} - (l+1)^{-1}\sum E_{ii}\) (Chap. VI, §4, no. 7.VIII), that is, consisting of the diagonal matrices of trace zero of the form \(x + (l+1)^{-1}a.1\) , where \(x\) has integer entries and \(a\in\mathbf{Z}\) . It follows that \(\mathfrak{sl}(l+1,\mathbf{Z})\) is the Chevalley order in \((\mathfrak{g},\mathfrak{h})\) associated to the permissible lattice \(\mathrm{Q}(\mathrm{R}^{\vee})\) and the Chevalley system \((X_{\alpha})\) . It is easy to verify that \(\bigwedge^{r}\mathbf{Z}^{l+1}\) is an admissible lattice in \(\bigwedge^{r}\mathbf{Q}^{l+1}\) relative to \(\mathfrak{sl}(l+1,\mathbf{Z})\) (§12, no. 8, Def. 3).

On the other hand, \(\mathfrak{gl}(l+1,\mathbf{Z})\) is a Chevalley order in the split reductive algebra \(\mathfrak{gl}(l+1,\mathbf{Q})\) ; its projection onto \(\mathfrak{sl}(l+1,\mathbf{Q})\) parallel to the centre \(\mathbf{Q}.1\) of \(\mathfrak{gl}(l+1,\mathbf{Q})\) is the Chevalley order in \((\mathfrak{g},\mathfrak{h})\) defined by the permissible lattice \(\mathrm{P}(\mathrm{R}^{\vee})\) in \(\mathfrak{h}\) and the Chevalley system \((X_{\alpha})\) . We remark that \(\mathfrak{gl}(l+1,\mathbf{Z})\) is not the direct sum of its intersections with \(\mathfrak{sl}(l+1,\mathbf{Q})\) and the centre of \(\mathfrak{gl}(l+1,\mathbf{Q})\) .